\documentclass[10pt]{report}
\usepackage{anyfontsize}
\AtBeginDocument{\fontsize{8.5}{10.6}\selectfont}

% ===================== ENGINE GUARD =====================
\usepackage{iftex}
\ifXeTeX\else
  \errmessage{This document must be compiled with XeLaTeX.}
\fi

% ===================== PAGE LAYOUT =====================
\usepackage[top=2cm,bottom=2.5cm,left=3cm,right=2cm]{geometry}
\setlength{\headheight}{15pt}

% ===================== CORE PACKAGES =====================
\usepackage{graphicx}
\usepackage{float}
\usepackage{array}
\usepackage{tabularx}
\usepackage{booktabs}
\usepackage{multirow}
\usepackage{threeparttable}
\usepackage{colortbl}
\usepackage{enumitem}
\usepackage{amsmath}
\usepackage{amssymb}
\usepackage{ragged2e}

% ===================== VISUAL DESIGN =====================
\usepackage[table]{xcolor}
\usepackage{tcolorbox}
\tcbuselibrary{breakable,skins,theorems}
\usepackage{framed}

\usepackage{fancyhdr}
\usepackage{titlesec}
\usepackage{tocloft}

% ===================== HEADER/FOOTER STYLING =====================
\pagestyle{fancy}
\fancyhf{}
\fancyhead[RE]{\nouppercase{\rightmark}}
\fancyhead[LO]{\nouppercase{\leftmark}}
\fancyhead[LE,RO]{\thepage}
\renewcommand{\headrulewidth}{0.4pt}
\renewcommand{\footrulewidth}{0pt}
\fancypagestyle{plain}{%
  \fancyhf{}
  \fancyfoot[C]{\thepage}
  \renewcommand{\headrulewidth}{0pt}
}

% ===================== TIKZ =====================
\usepackage{tikz}
\usetikzlibrary{shapes,arrows,positioning,calc,decorations.pathreplacing}

% ===================== PERSIAN (XeLaTeX) =====================
\usepackage{xepersian}

% ============ LATIN FONT SETUP ============
\setlatintextfont[
  Scale=0.95,
  Numbers=Lining,
  Ligatures=TeX
]{TeX Gyre Termes}

% ============ PERSIAN FONT SETUP - XB-KHORAMSHAHR ============
% ============ PERSIAN FONT SETUP - XB ZAR ============
\settextfont[
  Path=fonts/,
  Extension=.ttf,
  UprightFont=*,
  BoldFont=*Bd,
  ItalicFont=*It,
  BoldItalicFont=*BdIt,
  Script=Arabic,
  Language=Persian
]{XB-Yagut}  % تغییر از XB-Khoramshahr به XB Zar

\setdigitfont[
  Path=fonts/,
  Extension=.ttf,
  UprightFont=*,
  BoldFont=*Bd,
  ItalicFont=*It,
  BoldItalicFont=*BdIt,
  Numbers=Proportional,
  Script=Arabic,
  Language=Persian
]{XB-Yagut}  % تغییر از XB-Khoramshahr به XB Zar

\setmathdigitfont[
  Path=fonts/,
  Extension=.ttf,
  Script=Arabic,
  Language=Persian
]{XB-Khoramshahr}

\DefaultMathsDigits

% ============ CHAPTER/TITLE FORMATTING ============
\titleformat{\chapter}[display]
  {\normalfont\huge\bfseries\color{headerColor}}
  {\flushright\chaptertitlename\ \thechapter}
  {20pt}
  {\Huge\flushright}
\titlespacing*{\chapter}{0pt}{-20pt}{30pt}

\titleformat{\section}
  {\normalfont\Large\bfseries\color{headerColor}}
  {\thesection}{1em}{}

% ============ TABLE OF CONTENTS STYLING ============
\renewcommand{\cftchapfont}{\bfseries\color{headerColor}}
\renewcommand{\cftchapleader}{\cftdotfill{\cftdotsep}}
\setlength{\cftbeforechapskip}{6pt}
\renewcommand{\contentsname}{\hfill\Large\bfseries\color{headerColor}فهرست مطالب\hfill}

% ============ ENHANCED MACROS FOR BILINGUAL CONTENT ============
\newcommand{\lat}[1]{\lr{#1}}
\newcommand{\latinterm}[1]{\lr{\textit{#1}}}
\newcommand{\latinbold}[1]{\lr{\textbf{#1}}}
\newcommand{\latinnum}[1]{\lr{#1}}
\newcommand{\USD}[1]{\lr{\$#1}}
\newcommand{\percent}[1]{\lr{#1\%}}
\newcommand{\technical}[1]{\lr{\textsc{#1}}}

% ============ TABLE COLUMN TYPES FOR RTL ============
\newcolumntype{P}[1]{>{\RaggedRight}p{#1}}
\newcolumntype{H}[1]{>{\bfseries\color{headerColor}\RaggedRight}p{#1}}
\newcolumntype{C}[1]{>{\centering\arraybackslash}p{#1}}

% ============ LINE SPACING ============
\linespread{1.15}

% ============ HYPERREF SETUP ============
\usepackage{hyperref}
\hypersetup{
    colorlinks=true,
    linkcolor=headerColor,
    citecolor=headerColor,
    filecolor=headerColor,
    urlcolor=headerColor,
    pdftitle={مسیرهای ورود به Nasdaq Global Market: راهنمای جامع چهار مسیر استاندارد},
    pdfauthor={دکتر محسن دیرباز},
    pdfsubject={تحلیل مسیرهای پذیرش در بازار جهانی Nasdaq},
    pdfkeywords={Nasdaq, IPO, پذیرش اولیه, بازار سرمایه, استانداردهای پذیرش},
    bookmarksopen=true,
    pdfpagemode=UseOutlines
}

% ===================== PROFESSIONAL COLOR SCHEME =====================
\definecolor{headerColor}{RGB}{44,62,80}
\definecolor{accentColor}{RGB}{90, 70, 70}
\definecolor{pathwayIncome}{RGB}{235, 242, 248}
\definecolor{pathwayEquity}{RGB}{248, 243, 238}
\definecolor{pathwayMarket}{RGB}{238, 246, 240}
\definecolor{pathwayAssets}{RGB}{250, 248, 238}
\definecolor{conceptBox}{RGB}{240, 244, 248}
\definecolor{warningBox}{RGB}{250, 246, 240}
\definecolor{exampleBox}{RGB}{242, 248, 244}
\definecolor{summaryBox}{RGB}{248, 246, 242}
\definecolor{processBox}{RGB}{242, 248, 250}
\definecolor{requirementBox}{RGB}{245, 245, 250}
\definecolor{tableHeader}{RGB}{44, 62, 80}
\definecolor{tableRow1}{RGB}{255, 255, 255}
\definecolor{tableRow2}{RGB}{248, 249, 250}
\definecolor{frameBlue}{RGB}{70, 90, 120}
\definecolor{frameGreen}{RGB}{80, 110, 90}
\definecolor{frameOrange}{RGB}{150, 110, 70}
\definecolor{framePurple}{RGB}{100, 85, 120}
\definecolor{statusReady}{RGB}{80, 120, 90}
\definecolor{statusGap}{RGB}{150, 100, 70}
\definecolor{statusBlock}{RGB}{130, 70, 70}

% ===================== NAVIGATION SYSTEM =====================
\newcommand{\returnhome}{\vspace{0.5em}\noindent\protect\hyperlink{home}{\colorbox{conceptBox}{\color{headerColor}\textbf{بازگشت به فهرست}}}\par}
\newcommand{\returnscenarios}{\vspace{0.5em}\noindent\protect\hyperlink{scenarios}{\colorbox{conceptBox}{\color{headerColor}\textbf{بازگشت به سناریوها}}}\par}

% ===================== PROFESSIONAL TCOLORBOX STYLES =====================
\newtcolorbox{keyconceptbox}[1][]{
  enhanced,
  colback=conceptBox,
  colframe=headerColor,
  fonttitle=\bfseries,
  title={\color{white}مفهوم کلیدی},
  coltitle=white,
  attach boxed title to top right={yshift=-2mm, xshift=-4mm},
  boxed title style={colback=headerColor, sharp corners, size=small},
  breakable,
  sharp corners,
  boxrule=0.75pt,
  left=4mm,
  right=4mm,
  top=3mm,
  bottom=3mm,
  #1
}

\newtcolorbox{principlebox}[2][]{
  enhanced,
  colback=conceptBox,
  colframe=headerColor,
  fonttitle=\bfseries,
  title={#2},
  breakable,
  sharp corners,
  boxrule=0.75pt,
  left=4mm,
  right=4mm,
  top=3mm,
  bottom=3mm,
  #1
}

\newtcolorbox{warningbox}[2][]{
  enhanced,
  colback=warningBox,
  colframe=accentColor,
  fonttitle=\bfseries,
  title={#2},
  breakable,
  sharp corners,
  boxrule=0.75pt,
  left=4mm,
  right=4mm,
  top=3mm,
  bottom=3mm,
  #1
}

\newtcolorbox{examplebox}[2][]{
  enhanced,
  colback=exampleBox,
  colframe=frameGreen,
  fonttitle=\bfseries,
  title={#2},
  breakable,
  sharp corners,
  boxrule=0.75pt,
  left=4mm,
  right=4mm,
  top=3mm,
  bottom=3mm,
  #1
}

\newtcolorbox{processbox}[2][]{
  enhanced,
  colback=processBox,
  colframe=frameBlue,
  fonttitle=\bfseries,
  title={#2},
  breakable,
  sharp corners,
  boxrule=0.75pt,
  left=4mm,
  right=4mm,
  top=3mm,
  bottom=3mm,
  #1
}

\newtcolorbox{summarybox}[2][]{
  enhanced,
  colback=summaryBox,
  colframe=frameOrange,
  fonttitle=\bfseries,
  title={#2},
  breakable,
  sharp corners,
  boxrule=1pt,
  left=4mm,
  right=4mm,
  top=3mm,
  bottom=3mm,
  #1
}

\newtcolorbox{requirementbox}[2][]{
  enhanced,
  colback=requirementBox,
  colframe=framePurple,
  fonttitle=\bfseries,
  title={#2},
  breakable,
  sharp corners,
  boxrule=0.75pt,
  left=4mm,
  right=4mm,
  top=3mm,
  bottom=3mm,
  #1
}

\newtcolorbox{pathwayincomebox}[2][]{
  enhanced,
  colback=pathwayIncome,
  colframe=frameBlue,
  fonttitle=\bfseries,
  title={#2},
  breakable,
  sharp corners,
  boxrule=0.75pt,
  left=4mm,
  right=4mm,
  top=3mm,
  bottom=3mm,
  #1
}

\newtcolorbox{pathwayequitybox}[2][]{
  enhanced,
  colback=pathwayEquity,
  colframe=frameOrange,
  fonttitle=\bfseries,
  title={#2},
  breakable,
  sharp corners,
  boxrule=0.75pt,
  left=4mm,
  right=4mm,
  top=3mm,
  bottom=3mm,
  #1
}

\newtcolorbox{pathwaymarketbox}[2][]{
  enhanced,
  colback=pathwayMarket,
  colframe=frameGreen,
  fonttitle=\bfseries,
  title={#2},
  breakable,
  sharp corners,
  boxrule=0.75pt,
  left=4mm,
  right=4mm,
  top=3mm,
  bottom=3mm,
  #1
}

\newtcolorbox{pathwayassetsbox}[2][]{
  enhanced,
  colback=pathwayAssets,
  colframe=framePurple,
  fonttitle=\bfseries,
  title={#2},
  breakable,
  sharp corners,
  boxrule=0.75pt,
  left=4mm,
  right=4mm,
  top=3mm,
  bottom=3mm,
  #1
}

% ===================== TIKZ STYLES =====================
\tikzstyle{pathway1}=[rectangle, draw=frameBlue, fill=pathwayIncome, text width=3cm, text centered, rounded corners, minimum height=1.2cm, line width=1pt]
\tikzstyle{pathway2}=[rectangle, draw=frameOrange, fill=pathwayEquity, text width=3cm, text centered, rounded corners, minimum height=1.2cm, line width=1pt]
\tikzstyle{pathway3}=[rectangle, draw=frameGreen, fill=pathwayMarket, text width=3cm, text centered, rounded corners, minimum height=1.2cm, line width=1pt]
\tikzstyle{pathway4}=[rectangle, draw=framePurple, fill=pathwayAssets, text width=3cm, text centered, rounded corners, minimum height=1.2cm, line width=1pt]
\tikzstyle{arrow}=[->, >=stealth, line width=1.5pt]

% ===================== DOCUMENT STARTS =====================
\begin{document}

% ===================== TITLE PAGE =====================
\thispagestyle{empty}

\begin{center}
\vspace*{1.5cm}

{\Huge\bfseries\color{headerColor} مسیرهای ورود به \lat{Nasdaq Global Market}}

\vspace{0.6cm}

{\color{gray}\rule{0.7\textwidth}{0.75pt}}

\vspace{0.6cm}

{\Large\bfseries تحلیل چهار مسیر استاندارد برای شرکت‌های بین‌المللی}

\vspace{0.2cm}

{\large راهنمای جامع بر اساس استانداردهای \lat{Nasdaq Rule 5405}}

\vspace{1cm} % کاهش فاصله

{\LARGE\bfseries\color{headerColor} گردآورنده: دکتر محسن دیرباز}

\vspace{1cm} % کاهش فاصله

% in preamble (XeLaTeX/LuaLaTeX):
% \usepackage{bidi}

\begin{tikzpicture}[scale=1.2, transform shape]
    \node[pathway1, align=center] (p1) at (0,2)
      {\RL{\textbf{مسیر درآمد}}\\[-2pt]\small \LR{Income Standard}};

    \node[pathway2, align=center] (p2) at (4,2)
      {\RL{\textbf{مسیر سرمایه}}\\[-2pt]\small \LR{Equity Standard}};

    \node[pathway3, align=center] (p3) at (0,0)
      {\RL{\textbf{مسیر ارزش بازار}}\\[-2pt]\small \LR{Market Value Standard}};

    \node[pathway4, align=center] (p4) at (4,0)
      {\RL{\textbf{مسیر دارایی/درآمد}}\\[-2pt]\small \LR{Assets/Revenue Standard}};

    \node[draw=headerColor, fill=white, circle, minimum size=1.5cm, line width=1.5pt] at (2,1) {\Large\bfseries \lat{Nasdaq}};
\end{tikzpicture}

\vspace{2cm}

{\color{gray} نسخه ۱.۰ — تاریخ انتشار: ۱۰ بهمن ۱۴۰۴}

\vspace{0.2cm}

{\color{gray}\small (\lat{Version 1.0 — Publication Date: 30 January 2026})}

\vfill

{\small\color{gray} این سند بر اساس \lat{Nasdaq Rule 5405} و تجزیه‌وتحلیل بازار تهیه شده است.}\\
{\small\color{gray} کلیه ارقام به دلار آمریکا (\lat{USD}) بوده و مبنای محاسبات استانداردهای پذیرش اولیه است.}

\end{center}

% ===================== TABLE OF CONTENTS =====================
\clearpage
\hypertarget{home}{}
\tableofcontents

% ===================== INTRODUCTION =====================
\clearpage
\chapter*{مقدمه}
\addcontentsline{toc}{chapter}{مقدمه}

\begin{keyconceptbox}
\lat{Nasdaq Global Market} چهار مسیر مجزا برای پذیرش اولیه ارائه می‌دهد که هر یک بر اساس نقاط قوت مالی و عملیاتی متفاوت شرکت‌ها طراحی شده‌اند. این سند با رویکردی تحلیلی، هر مسیر را تشریح کرده و با سناریوهای واقعی نشان می‌دهد چگونه شرکت‌های مختلف می‌توانند حداقل الزامات را برآورده کنند.
\end{keyconceptbox}

\section*{ساختار و روش‌شناسی سند}

این راهنما بر اساس چارچوب سه‌بخشی زیر طراحی شده است:

\begin{enumerate}[rightmargin=1.5cm]
\item \textbf{الزامات پایه}: شروط مشترک و غیرقابل مذاکره برای همه شرکت‌های متقاضی
\item \textbf{چهار مسیر متمایز}: تشریح کامل هر استاندارد پذیرش با تحلیل منطق طراحی
\item \textbf{سناریوهای عملی}: نقشه‌برداری فاصله تا استاندارد برای حالت‌های واقعی بازار
\end{enumerate}

\section*{رویکرد تحلیلی و محدوده بررسی}

\begin{principlebox}{مبانی تحلیلی}
این تحلیل بر سه اصل استوار است:
\begin{itemize}[rightmargin=1.5cm]
\item \textbf{واقع‌گرایی مالی}: تمام سناریوها بر اساس داده‌های واقعی بازار طراحی شده‌اند
\item \textbf{شفافیت معیارها}: هر آستانه عددی با منطق طراحی آن توضیح داده شده است
\item \textbf{قابلیت اجرا}: همه راه‌حل‌های پیشنهادی در چارچوب عملیاتی امکان‌پذیر هستند
\end{itemize}
\end{principlebox}

\textbf{محدوده بررسی:} این سند فقط استانداردهای پذیرش اولیه (\lat{Initial Listing Standards}) را پوشش می‌دهد. استانداردهای تداوم حضور (\lat{Continued Listing Standards}) و سایر مقررات \lat{Nasdaq} نیازمند بررسی جداگانه هستند.

% ===================== CHAPTER 1: BASE REQUIREMENTS =====================
\clearpage
\chapter{الزامات پایه برای همه مسیرها}
\hypertarget{base}{}

\begin{requirementbox}{شروط غیرقابل مذاکره}
صرف‌نظر از مسیر انتخابی، همه شرکت‌های متقاضی باید این شروط اساسی را به‌طور کامل برآورده کنند. این الزامات تضمین‌کننده حداقل استانداردهای بازار برای سرمایه‌گذاران هستند.
\end{requirementbox}

\section{الزامات ساختاری و عملیاتی}

\begin{table}[H]
\centering
\caption{الزامات پایه \lat{Nasdaq Global Market} — مشترک در همه مسیرها}
\label{tab:base_requirements}
\begin{threeparttable}
\begin{tabularx}{\textwidth}{H{4.5cm} X C{2.5cm}}
\toprule
\rowcolor{tableHeader}
\textcolor{white}{\textbf{الزام}} & \textcolor{white}{\textbf{شرح و منطق طراحی}} & \textcolor{white}{\textbf{آستانه}} \\
\midrule
\rowcolor{conceptBox}
قیمت پیشنهادی اولیه & حداقل قیمت سهم در \lat{IPO} برای اطمینان از نقدشوندگی اولیه & \USD{4} >= \\
\rowcolor{tableRow2}
سهام عمومی قابل معامله & تعداد سهام قابل معامله آزاد برای تضمین عمق بازار & \lat{1.1M} سهم \\
\rowcolor{requirementBox}
سهامداران مختلف & تعداد افراد/نهادهای مجزا برای ایجاد پایه سهامداران گسترده & \lat{400} نفر\tnote{a} \\
\rowcolor{tableRow2}
صورت‌های مالی حسابرسی‌شده & سوابق مالی معتبر برای شفافیت و قابلیت اتکا & ۳ سال متوالی\tnote{b} \\
\rowcolor{conceptBox}
بازارسازهای متعهد & نهادهای تضمین‌کننده نقدشوندگی در روزهای اولیه معاملات & \lat{3} یا \lat{4}\tnote{c} \\
\bottomrule
\end{tabularx}
\begin{tablenotes}\small
\item[a] حداقل ۵۰٪ این سهامداران باید دارایی \USD{2,500} یا بیشتر داشته باشند
\item[b] طبق استانداردهای \lat{GAAP} یا \lat{IFRS} با تأیید حسابرس مستقل
\item[c] مسیرهای درآمدی و سرمایه: \lat{3} بازارساز — مسیرهای بازار و دارایی: \lat{4} بازارساز
\end{tablenotes}
\end{threeparttable}
\end{table}
\clearpage
\section{یادداشت‌های تحلیلی و هشدارهای عملی}

\begin{warningbox}{نکات حیاتی اجرایی}
\begin{itemize}[rightmargin=1.5cm]
\item \textbf{قیمت \USD{4}}: این قیمت پیشنهادی اولیه در \lat{IPO} است، نه ارزش بازار روزانه. شرکت باید قیمت‌گذاری \lat{IPO} را با درنظرگرفتن این محدودیت برنامه‌ریزی کند.
\item \textbf{سهام عمومی}: ۱٫۱ میلیون سهم باید بدون محدودیت قفل‌شدگی (\lat{lock-up}) و قابل معامله آزاد باشد.
\item \textbf{سهامداران}: تعداد واقعی افراد/نهادهای مستقل ملاک است، نه تعداد حساب‌های کارگزاری.
\item \textbf{حسابرسی}: باید توسط حسابرسان مستقل واجد شرایط \lat{PCAOB} انجام شده و گزارش حسابرسی بدون شرط (\lat{unqualified opinion}) ارائه شود.
\end{itemize}
\end{warningbox}

\returnhome

% ===================== CHAPTER 2: THE FOUR PATHWAYS =====================
\clearpage
\chapter{چهار مسیر متمایز پذیرش}
\hypertarget{pathways}{}

\begin{keyconceptbox}
هر مسیر پذیرش بر اساس نوع متفاوتی از قدرت شرکت طراحی شده است. انتخاب مسیر باید مبتنی بر نقاط قوت فعلی و برنامه‌های توسعۀ شرکت باشد، نه آرزوها یا projections غیرواقعی.
\end{keyconceptbox}

\section{نمای کلی و منطق طراحی مسیرها}

\begin{table}[H]
\centering
\caption{مقایسه استراتژیک چهار مسیر پذیرش}
\label{tab:pathways_overview}
\begin{threeparttable}
\begin{tabularx}{\textwidth}{H{3cm} X P{2.8cm} P{2.8cm}}
\toprule
\rowcolor{tableHeader}
\textcolor{white}{\textbf{مسیر}} & \textcolor{white}{\textbf{منطق طراحی و فلسفه}} & \textcolor{white}{\textbf{شرکت‌های مناسب}} & \textcolor{white}{\textbf{چالش متداول}} \\
\midrule
\rowcolor{pathwayIncome}
۱. درآمدی & سودآوری اثبات‌شده + سرمایه متوسط = ریسک پایین برای بازار & شرکت‌های بالغ سودده با رشد پایدار & حقوق سهام محدود \\
\rowcolor{pathwayEquity}
۲. سرمایه & سرمایه قوی + سابقه عملیاتی = بلوغ سازمانی & استارتاپ‌های سرمایه‌دار در مراحل رشد & الزام زمانی سخت (۲ سال) \\
\rowcolor{pathwayMarket}
۳. ارزش بازار & ارزش‌گذاری بالا توسط بازار = اعتبار و مقیاس & شرکت‌های پررشد (یونیکورن‌ها) & وابستگی به شرایط بازار \\
\rowcolor{pathwayAssets}
۴. دارایی/درآمد & مقیاس عملیاتی بالا = پایداری و تاب‌آوری & شرکت‌های تولیدی، زیرساختی، املاک & نیاز به حجم عملیات بالا \\
\bottomrule
\end{tabularx}
\end{threeparttable}
\end{table}
\clearpage

\section{مسیر ۱: استاندارد درآمدی (\lat{Income Standard})}

\begin{pathwayincomebox}{مسیر درآمدی: اثبات سودآوری}
این مسیر برای شرکت‌هایی طراحی شده که توانایی کسب سود پایدار را اثبات کرده‌اند، حتی اگر حقوق سهام آن‌ها نسبتاً متوسط باشد. منطق اصلی: «سودآوری واقعی، ریسک پایین‌تر».
\end{pathwayincomebox}

\subsection{الزامات کمی و کیفی}

\begin{table}[H]
\centering
\caption{الزامات دقیق مسیر درآمدی — \lat{Income Standard}}
\label{tab:income_requirements}
\begin{threeparttable}
\begin{tabularx}{\textwidth}{H{5cm} X C{2.5cm}}
\toprule
\rowcolor{tableHeader}
\textcolor{white}{\textbf{الزام}} & \textcolor{white}{\textbf{تعریف عملیاتی و دامنه شمول}} & \textcolor{white}{\textbf{آستانه}} \\
\midrule
\rowcolor{pathwayIncome}
حقوق سهام (\lat{Equity}) & کل دارایی‌ها منهای کل بدهی‌ها در ترازنامه & \USD{15M} >= \\
\rowcolor{tableRow2}
سود عملیاتی قبل از مالیات & درآمد حاصل از عملیات اصلی و ادامه‌دار شرکت & \USD{1M}\tnote{a} \\
\rowcolor{pathwayIncome}
ارزش بازار سهام عمومی & ارزش کل سهام قابل معامله آزاد در قیمت \lat{IPO} & \USD{15M} >= \\
\rowcolor{tableRow2}
بازارسازهای متعهد & نهادهای تضمین‌کننده نقدشوندگی در ۳ ماه اول & \lat{3} >= \\
\bottomrule
\end{tabularx}
\begin{tablenotes}\small
\item[a] در سال مالی اخیر یا در ۲ سال از ۳ سال گذشته — مبنای حسابداری: \lat{GAAP/IFRS}
\end{tablenotes}
\end{threeparttable}
\end{table}

\subsection{تحلیل منطق و مبادلات طراحی}

\begin{principlebox}{چرا این ترکیب خاص؟}
این مسیر مبتنی بر مبادله استراتژیک زیر طراحی شده است:

\textbf{مبادله طراحی:} حقوق سهام پایین‌تر (\USD{15M} در مقابل \USD{30M} در مسیر سرمایه) در مقابل سودآوری اثبات‌شده (\USD{1M}+)

\textbf{منطق بازار:} شرکت‌های سودآور حتی با سرمایه متوسط، ریسک کمتری برای سرمایه‌گذاران خرد ایجاد می‌کنند. سودآوری نشان‌دهنده مدل کسب‌وکار کارا، مدیریت اثربخش و جریان نقدی مثبت است.
\end{principlebox}
\clearpage

\section{مسیر ۲: استاندارد سرمایه (\lat{Equity Standard})}

\begin{pathwayequitybox}{مسیر سرمایه: قدرت ترازنامه}
این مسیر برای شرکت‌هایی است که سرمایه قوی جذب کرده‌اند ولی هنوز به سودآوری نرسیده‌اند. منطق اصلی: «سرمایه کافی، زمان برای رسیدن به سودآوری».
\end{pathwayequitybox}

\subsection{الزامات کمی و کیفی}

\begin{table}[H]
\centering
\caption{الزامات دقیق مسیر سرمایه — \lat{Equity Standard}}
\label{tab:equity_requirements}
\begin{threeparttable}
\begin{tabularx}{\textwidth}{H{5cm} X C{2.5cm}}
\toprule
\rowcolor{tableHeader}
\textcolor{white}{\textbf{الزام}} & \textcolor{white}{\textbf{تعریف عملیاتی و دامنه شمول}} & \textcolor{white}{\textbf{آستانه}} \\
\midrule
\rowcolor{pathwayEquity}
حقوق سهام (\lat{Equity}) & کل دارایی‌ها منهای کل بدهی‌ها — نشان‌دهنده قدرت مالی & \USD{30M} >= \\
\rowcolor{tableRow2}
سابقه عملیاتی مستمر & مدت زمان فعالیت شرکت به‌عنوان نهاد قانونی فعال & ۲ سال کامل \\
\rowcolor{pathwayEquity}
ارزش بازار سهام عمومی & ارزش کل سهام قابل معامله آزاد در قیمت \lat{IPO} & \USD{18M} >= \\
\rowcolor{tableRow2}
بازارسازهای متعهد & نهادهای تضمین‌کننده نقدشوندگی در ۳ ماه اول & \lat{3} >= \\
\bottomrule
\end{tabularx}
\end{threeparttable}
\end{table}

\subsection{تحلیل منطق و محدودیت‌های کلیدی}

\begin{principlebox}{منطق طراحی و هشدار اجرایی}
این مسیر مبتنی بر این فرض طراحی شده است که شرکت‌های با سرمایه کافی (\USD{30M}+) و سابقه عملیاتی (۲ سال) حتی بدون سودآوری، به اندازه کافی بالغ و کم‌ریسک هستند.

\textbf{مبادله طراحی:} حذف نیاز به سودآوری در مقابل الزام حقوق سهام دوبرابر (\USD{30M} در مقابل \USD{15M}) + الزام زمانی سخت (۲ سال)

\textbf{هشدار حیاتی:} ۲ سال سابقه عملیاتی یک مانع زمانی غیرقابل کاهش است. این مدت باید از تاریخ ثبت شرکت به‌عنوان یک نهاد قانونی فعال محاسبه شود.
\end{principlebox}

\clearpage
\section{مسیر ۳: استاندارد ارزش بازار (\lat{Market Value Standard})}

\begin{pathwaymarketbox}{مسیر ارزش بازار: تأیید اعتبار توسط بازار}
این مسیر برای شرکت‌های بسیار بزرگ طراحی شده که بازار سرمایه اعتبار آن‌ها را از طریق ارزش‌گذاری بالا تأیید کرده است. منطق اصلی: «مقیاس و اعتبار، جایگزین سودآوری».
\end{pathwaymarketbox}

\subsection{الزامات کمی و کیفی}

\begin{table}[H]
\centering
\caption{الزامات دقیق مسیر ارزش بازار — \lat{Market Value Standard}}
\label{tab:market_requirements}
\begin{threeparttable}
\begin{tabularx}{\textwidth}{H{5cm} X C{2.5cm}}
\toprule
\rowcolor{tableHeader}
\textcolor{white}{\textbf{الزام}} & \textcolor{white}{\textbf{تعریف عملیاتی و دامنه شمول}} & \textcolor{white}{\textbf{آستانه}} \\
\midrule
\rowcolor{pathwayMarket}
ارزش بازار کل اوراق (\lat{MVLS}) & ارزش کل تمام سهام و اوراق بدهی قابل معامله شرکت & \USD{75M} >= \\
\rowcolor{tableRow2}
ارزش بازار سهام عمومی & ارزش کل سهام قابل معامله آزاد در قیمت \lat{IPO} & \USD{20M} >= \\
\rowcolor{pathwayMarket}
بازارسازهای متعهد & نهادهای تضمین‌کننده نقدشوندگی در ۳ ماه اول & \lat{4} >= \\
\rowcolor{tableRow2}
شرط ویژه برای شرکت‌های موجود\tnote{a} & سازوکار تأیید اعتبار بازار برای شرکت‌های فهرست‌شده در سایر بورس‌ها & ۹۰ روز معامله \\
\bottomrule
\end{tabularx}
\begin{tablenotes}\small
\item[a] برای شرکت‌هایی که فقط از این مسیر استفاده می‌کنند: باید ۹۰ روز متوالی با \lat{MVLS} \USD{75M}+ و قیمت پیشنهاد \USD{4}+ در بازار قبلی معامله شده باشند
\end{tablenotes}
\end{threeparttable}
\end{table}

\subsection{تحلیل منطق و شرایط کاربردی}

\begin{principlebox}{منطق طراحی و کاربران هدف}
این مسیر مبتنی بر این اصل است که شرکت‌های با ارزش بازار بسیار بالا (\USD{75M}+) به‌اندازه‌ای بزرگ و شناخته‌شده هستند که الزامات سنتی سودآوری و سرمایه برای آن‌ها ضروری نیست.

\textbf{مبادله طراحی:} حذف تمام الزامات سودآوری و حقوق سهام در مقابل ارزش‌گذاری بسیار بالا

\textbf{کاربران هدف معمول:}
\begin{itemize}[rightmargin=1.5cm]
\item یونیکورن‌ها (ارزش‌گذاری \USD{1B}+)
\item شرکت‌های پررشد با سرمایه‌گذاری سنگین در فاز توسعه
\item شرکت‌های فناوری با مقیاس بزرگ اما زیان‌ده در کوتاه‌مدت
\end{itemize}
\end{principlebox}
\clearpage
\section{مسیر ۴: استاندارد دارایی/درآمد (\lat{Total Assets/Revenue Standard})}

\begin{pathwayassetsbox}{مسیر دارایی/درآمد: مقیاس عملیاتی اثبات‌شده}
این مسیر برای شرکت‌های بزرگ با عملیات سنگین دارایی طراحی شده است. منطق اصلی: «مقیاس عملیاتی بزرگ، نشان‌دهنده پایداری و تاب‌آوری».
\end{pathwayassetsbox}

\subsection{الزامات کمی و کیفی}

\begin{table}[H]
\centering
\caption{الزامات دقیق مسیر دارایی/درآمد — \lat{Assets/Revenue Standard}}
\label{tab:assets_requirements}
\begin{threeparttable}
\begin{tabularx}{\textwidth}{H{5cm} X C{2.5cm}}
\toprule
\rowcolor{tableHeader}
\textcolor{white}{\textbf{الزام}} & \textcolor{white}{\textbf{تعریف عملیاتی و دامنه شمول}} & \textcolor{white}{\textbf{آستانه}} \\
\midrule
\rowcolor{pathwayAssets}
مجموع دارایی‌های شرکت & کل دارایی‌های مشهود و نامشهود در ترازنامه & \USD{75M} >= \\
\rowcolor{tableRow2}
مجموع درآمدهای عملیاتی & کل فروش و درآمدهای عملیاتی سالانه شرکت & \USD{75M}\tnote{a} \\
\rowcolor{pathwayAssets}
ارزش بازار سهام عمومی & ارزش کل سهام قابل معامله آزاد در قیمت \lat{IPO} & \USD{20M} >= \\
\rowcolor{tableRow2}
بازارسازهای متعهد & نهادهای تضمین‌کننده نقدشوندگی در ۳ ماه اول & \lat{4} >= \\
\bottomrule
\end{tabularx}
\begin{tablenotes}\small
\item[a] در سال مالی اخیر یا در ۲ سال از ۳ سال گذشته — درآمد باید از عملیات اصلی شرکت باشد
\end{tablenotes}
\end{threeparttable}
\end{table}

\subsection{تحلیل منطق و صنایع هدف}

\begin{principlebox}{منطق طراحی و مصادیق عملی}
این مسیر مبتنی بر این اصل است که شرکت‌های با دارایی و درآمد بالا (\USD{75M} هر کدام) به‌اندازه‌ای بزرگ هستند که مقیاس عملیاتی آن‌ها خود دلیلی بر بلوغ و پایداری است.

\textbf{مبادله طراحی:} حذف الزامات سودآوری و حقوق سهام در مقابل دارایی بالا \textbf{و} درآمد بالا (شرایط ترکیبی)

\textbf{صنایع هدف اصلی:}
\begin{itemize}[rightmargin=1.5cm]
\item تولیدی و صنعتی (تجهیزات، ماشین‌آلات)
\item املاک و مستغلات (توسعه‌دهندگان، مالکان)
\item زیرساخت و انرژی (پروژه‌های بزرگ)
\item خرده‌فروشی زنجیره‌ای با دارایی فیزیکی زیاد
\end{itemize}
\end{principlebox}

\returnhome

% ===================== CHAPTER 3: PRACTICAL SCENARIOS =====================
\clearpage
\chapter{سناریوهای عملی: تحلیل فاصله تا استاندارد}
\hypertarget{scenarios}{}

\begin{keyconceptbox}
این فصل چهار سناریوی واقعی ارائه می‌دهد که نشان می‌دهد شرکت‌های مختلف چگونه می‌توانند به حداقل الزامات هر مسیر برسند. هر سناریو از یک نقطه شروع منطقی آغاز می‌شود و دقیقاً مشخص می‌کند چه فاصله‌هایی باید پر شود.
\end{keyconceptbox}

\section{اصول طراحی سناریو و مفروضات تحلیلی}

\begin{principlebox}{چارچوب طراحی سناریوها}
هر سناریو بر اساس اصول زیر طراحی شده است:
\begin{itemize}[rightmargin=1.5cm]
\item \textbf{نقطه شروع واقع‌بینانه}: هر سناریو با موقعیتی معقول و مبتنی بر الگوهای واقعی بازار آغاز می‌شود
\item \textbf{هدف حداقلی}: تمرکز بر رسیدن به دقیقاً آستانه‌های الزامی است، نه فراتر از آن
\item \textbf{فاصله‌های کمّی}: هر فاصله به صورت عددی و قابل اندازه‌گیری بیان شده است
\item \textbf{زمان‌بندی واقع‌بینانه}: برآوردهای زمانی مبتنی بر تجربه صنعت و محدودیت‌های عملیاتی
\end{itemize}
\end{principlebox}

% SCENARIO 1
\clearpage
\section{سناریو ۱: شرکت \lat{SaaS} در مرز سودآوری}

\begin{examplebox}{پروفایل شرکت و موقعیت بازار}
\textbf{نوع کسب‌وکار:} پلتفرم \lat{SaaS B2B} برای مدیریت زنجیره تأمین\\
\textbf{سن شرکت:} ۴ سال از تأسیس\\
\textbf{موقعیت مالی:} تازه به نقطه سر به سر (\lat{breakeven}) رسیده، درآمد رو به رشد\\
\textbf{مسیر هدف:} استاندارد درآمدی (\lat{Income Standard})\\
\textbf{منطق انتخاب مسیر:} سودآوری نوپا اما واقعی، حقوق سهام محدود اما در حال رشد
\end{examplebox}

\subsection{موقعیت فعلی (نقطه شروع تحلیلی)}

\begin{table}[H]
\centering
\caption{سناریو ۱ — وضعیت مالی و عملیاتی فعلی}
\label{tab:scenario1_current}
\begin{threeparttable}
\begin{tabularx}{\textwidth}{H{4.5cm} X C{2.5cm}}
\toprule
\rowcolor{tableHeader}
\textcolor{white}{\textbf{شاخص کلیدی}} & \textcolor{white}{\textbf{توضیح تحلیلی و اهمیت استراتژیک}} & \textcolor{white}{\textbf{مقدار فعلی}} \\
\midrule
\rowcolor{tableRow1}
درآمد سالانه تکراری (\lat{ARR}) & درآمد مکرر سالانه — نشان‌دهنده پایداری درآمد & \USD{8.5M} \\
\rowcolor{conceptBox}
سود قبل از بهره، مالیات و استهلاک (\lat{EBITDA}) & سود عملیاتی — نشان‌دهنده کارایی عملیاتی & \USD{425K}\tnote{a} \\
\rowcolor{tableRow1}
حقوق سهام (\lat{Equity}) & کل دارایی‌ها منهای بدهی‌ها — قدرت ترازنامه & \USD{9M} \\
\rowcolor{requirementBox}
تأمین مالی تجمعی & کل سرمایه جذب‌شده در دورهای \lat{Seed} و \lat{Series A} & \USD{12M} \\
\rowcolor{tableRow1}
ارزش‌گذاری آخرین دور & ارزش‌گذاری در دور \lat{Series A} (۶ ماه پیش) & \USD{35M} \\
\rowcolor{conceptBox}
نرخ رشد سالانه درآمد (\lat{CAGR}) & میانگین رشد مرکب سالانه در ۲ سال گذشته & ۶۵٪ \\
\bottomrule
\end{tabularx}
\begin{tablenotes}\small
\item[a] حاشیه سود \lat{EBITDA}: ۵٪ — در محدوده معمول برای شرکت‌های \lat{SaaS} در حال رشد
\end{tablenotes}
\end{threeparttable}
\end{table}

\subsection{تحلیل فاصله تا استاندارد درآمدی}

\begin{table}[H]
\centering
\caption{سناریو ۱ — تحلیل فاصله کمی تا الزامات مسیر درآمدی}
\label{tab:scenario1_gap}
\begin{threeparttable}
\begin{tabularx}{\textwidth}{H{4cm} C{2.2cm} C{2.2cm} X}
\toprule
\rowcolor{tableHeader}
\textcolor{white}{\textbf{الزام}} & \textcolor{white}{\textbf{آستانه}} & \textcolor{white}{\textbf{وضعیت فعلی}} & \textcolor{white}{\textbf{تحلیل فاصله و اولویت}} \\
\midrule
\rowcolor{pathwayIncome}
حقوق سهام & \USD{15M} >= & \USD{9M} & \textcolor{statusGap}{\textbf{فاصله: \USD{6M}} — اولویت بالا: نیاز به تأمین مالی اضافی} \\
\rowcolor{tableRow2}
سود عملیاتی قبل از مالیات & \USD{1M} >= & \USD{425K} & \textcolor{statusGap}{\textbf{فاصله: \USD{575K}} — اولویت متوسط: نیاز به رشد درآمد یا بهبود حاشیه} \\
\rowcolor{pathwayIncome}
ارزش بازار سهام عمومی & \USD{15M} >= & \USD{18M}\tnote{a} & \textbf{برآورده شده} — فاصله: صفر \\
\rowcolor{tableRow2}
بازارسازهای متعهد & \lat{3} >= & تأمین خودکار\tnote{b} & \textbf{برآورده شده} — فاصله: صفر \\
\bottomrule
\end{tabularx}
\begin{tablenotes}\small
\item[a] فرض تحلیلی: \lat{float} ۵۰٪ در \lat{IPO} با ارزش‌گذاری \USD{35M} → \USD{17.5M} ارزش بازار سهام عمومی
\item[b] در \lat{IPO} استاندارد، بازارسازها توسط بانک سرمایه‌گذاری (\lat{underwriter}) تأمین می‌شوند
\end{tablenotes}
\end{threeparttable}
\end{table}

\subsection{مسیر عملیاتی رسیدن به آمادگی}

\begin{processbox}{برنامه اجرایی ۱۸ ماهه — فاز‌بندی شده}
\textbf{فاز ۱: رشد به سودآوری هدف (ماه‌های ۱–۹)}
\begin{itemize}[rightmargin=1.5cm]
\item \textbf{هدف درآمدی}: رسیدن به \lat{ARR} = \USD{12M} (رشد ۴۱٪ از سطح فعلی)
\item \textbf{هدف سودآوری}: افزایش حاشیه \lat{EBITDA} به ۱۰٪ → سود سالانه \USD{1.2M}
\item \textbf{اقدامات کلیدی اجرایی}:
  \begin{enumerate}[rightmargin=1.5cm]
  \item بهینه‌سازی هزینه کسب مشتری (\lat{CAC}) از طریق کانال‌های بازاریابی کاراتر
  \item کاهش نرخ ریزش مشتریان (\lat{Churn Rate}) به کمتر از ۵٪ سالانه
  \item افزایش درآمد به ازای هر کاربر (\lat{ARPU}) از طریق فروش اضافی (\lat{upselling})
  \end{enumerate}
\end{itemize}

\textbf{فاز ۲: تأمین مالی استراتژیک (ماه‌های ۱۰–۱۲)}
\begin{itemize}[rightmargin=1.5cm]
\item \textbf{هدف مالی}: جذب \USD{8M} در دور \lat{Series B}
\item \textbf{تخصیص منابع}: \USD{6M} برای تقویت ترازنامه (پر کردن فاصله حقوق سهام)، \USD{2M} برای سرمایه رشد
\item \textbf{ارزش‌گذاری هدف}: \USD{50–60M} (ضریب ۵ برابر درآمد سالانه برای \lat{SaaS})
\end{itemize}

\textbf{فاز ۳: آماده‌سازی \lat{IPO} (ماه‌های ۱۳–۱۸)}
\begin{itemize}[rightmargin=1.5cm]
\item تکمیل حسابرسی‌های مالی ۳ ساله مطابق \lat{GAAP}
\item ساختاردهی حاکمیت شرکتی (\lat{corporate governance}) با هیئت‌مدیره مستقل
\item انتخاب بانک سرمایه‌گذاری (\lat{underwriter}) و ثبت فرم \lat{S-1} نزد \lat{SEC}
\end{itemize}
\end{processbox}

\subsection{خلاصه تحلیلی و مدیریت ریسک}

\begin{summarybox}{نتیجه‌گیری سناریو و چارچوب زمانی}
\begin{tabularx}{\textwidth}{>{\bfseries\color{headerColor}}l X C{2cm}}
\textbf{فاصله اصلی شناسایی‌شده} & حقوق سهام (\USD{6M}) + سود عملیاتی (\USD{575K}) & — \\[4pt]
\textbf{راه‌حل عملیاتی پیشنهادی} & ترکیب رشد درآمد (فاز ۱) + تأمین مالی \USD{8M} (فاز ۲) & — \\[4pt]
\textbf{چارچوب زمانی کل} & ۱۸ ماه (۹ ماه رشد + ۳ ماه تأمین مالی + ۶ ماه آماده‌سازی \lat{IPO}) & \textbf{۱۸ ماه} \\[4pt]
\textbf{ریسک‌های اصلی اجرایی} & عدم دستیابی به اهداف رشد درآمد، کاهش تقاضای بازار برای \lat{SaaS} & سطح متوسط \\[4pt]
\textbf{شاخص‌های موفقیت کلیدی (\lat{KPIs})} & \lat{ARR} ≥ \USD{12M}، حاشیه \lat{EBITDA} ≥ ۱۰٪، جذب سرمایه \USD{8M}+ & — \\
\end{tabularx}
\end{summarybox}

\returnscenarios

% SCENARIO 2
\clearpage
\section{سناریو ۲: استارتاپ بیوتکنولوژی با سرمایه قوی}

\begin{examplebox}{پروفایل شرکت و موقعیت بازار}
\textbf{نوع کسب‌وکار:} شرکت داروسازی در مرحله کارآزمایی بالینی\\
\textbf{سن شرکت:} ۹ ماه از تأسیس رسمی\\
\textbf{موقعیت محصول:} در مرحله \lat{Phase 2 clinical trial}، بدون درآمد عملیاتی\\
\textbf{مسیر هدف:} استاندارد سرمایه (\lat{Equity Standard})\\
\textbf{منطق انتخاب مسیر:} سرمایه قوی، سابقه عملیاتی محدود، عدم سودآوری
\end{examplebox}

\subsection{موقعیت فعلی (نقطه شروع تحلیلی)}

\begin{table}[H]
\centering
\caption{سناریو ۲ — وضعیت مالی و عملیاتی فعلی}
\label{tab:scenario2_current}
\begin{threeparttable}
\begin{tabularx}{\textwidth}{H{4.5cm} X C{2.5cm}}
\toprule
\rowcolor{tableHeader}
\textcolor{white}{\textbf{شاخص کلیدی}} & \textcolor{white}{\textbf{توضیح تحلیلی و اهمیت استراتژیک}} & \textcolor{white}{\textbf{مقدار فعلی}} \\
\midrule
\rowcolor{tableRow1}
تأمین مالی تجمعی & کل سرمایه جذب‌شده در دورهای \lat{Seed} و \lat{Series A} & \USD{35M} \\
\rowcolor{conceptBox}
حقوق سهام (\lat{Equity}) & پس از کسر هزینه‌های عملیاتی (\lat{burn rate}) & \USD{32M}\tnote{a} \\
\rowcolor{tableRow1}
سابقه عملیاتی مستمر & مدت زمان فعالیت شرکت به‌عنوان نهاد قانونی فعال & ۹ ماه \\
\rowcolor{requirementBox}
ارزش‌گذاری آخرین دور & ارزش‌گذاری در دور \lat{Series A} (۳ ماه پیش) & \USD{120M} \\
\rowcolor{tableRow1}
وضعیت توسعه محصول & مرحله کارآزمایی بالینی داروی اصلی & \lat{Phase 2} \\
\rowcolor{conceptBox}
نرخ مصرف ماهانه (\lat{Monthly Burn Rate}) & متوسط هزینه‌های عملیاتی ماهانه & \USD{333K} \\
\bottomrule
\end{tabularx}
\begin{tablenotes}\small
\item[a] \USD{35M} جذب شده منهای \USD{3M} هزینه‌های عملیاتی در ۹ ماه گذشته
\end{tablenotes}
\end{threeparttable}
\end{table}

\subsection{تحلیل فاصله تا استاندارد سرمایه}

\begin{table}[H]
\centering
\caption{سناریو ۲ — تحلیل فاصله کمی تا الزامات مسیر سرمایه}
\label{tab:scenario2_gap}
\begin{threeparttable}
\begin{tabularx}{\textwidth}{H{4cm} C{2.2cm} C{2.2cm} X}
\toprule
\rowcolor{tableHeader}
\textcolor{white}{\textbf{الزام}} & \textcolor{white}{\textbf{آستانه}} & \textcolor{white}{\textbf{وضعیت فعلی}} & \textcolor{white}{\textbf{تحلیل فاصله و اولویت}} \\
\midrule
\rowcolor{pathwayEquity}
حقوق سهام & \USD{30M} >= & \USD{32M} & \textbf{برآورده شده} — فاصله: صفر \\
\rowcolor{tableRow2}
سابقه عملیاتی مستمر & ۲ سال >= & ۹ ماه & \textcolor{statusBlock}{\textbf{فاصله: ۱۵ ماه}} — اولویت حیاتی: مانع زمانی غیرقابل کاهش \\
\rowcolor{pathwayEquity}
ارزش بازار سهام عمومی & \USD{18M} >= & \USD{60M}\tnote{a} & \textbf{برآورده شده} — فاصله: صفر \\
\rowcolor{tableRow2}
بازارسازهای متعهد & \lat{3} >= & تأمین خودکار\tnote{b} & \textbf{برآورده شده} — فاصله: صفر \\
\bottomrule
\end{tabularx}
\begin{tablenotes}\small
\item[a] فرض تحلیلی: \lat{float} ۵۰٪ در \lat{IPO} با ارزش‌گذاری \USD{120M}
\item[b] در \lat{IPO} استاندارد، بازارسازها توسط بانک سرمایه‌گذاری تأمین می‌شوند
\end{tablenotes}
\end{threeparttable}
\end{table}

\subsection{مسیر عملیاتی و محدودیت زمانی}

\begin{warningbox}{هشدار حیاتی: مانع زمانی غیرقابل کاهش}
\textbf{محدودیت ساختاری:} الزام ۲ سال سابقه عملیاتی یک شرط زمانی سخت (\lat{hard stop}) است. هیچ مکانیزمی برای تسریع یا دور زدن این الزام وجود ندارد.

\textbf{زمان باقی‌مانده اجباری:} ۱۵ ماه تا رسیدن به حداقل ۲ سال سابقه عملیاتی

\textbf{تأثیر استراتژیک:} این محدودیت، زمان‌بندی \lat{IPO} را مستقل از عملکرد مالی یا رشد شرکت تعیین می‌کند.
\end{warningbox}

\begin{processbox}{برنامه اجرایی ۲۴ ماهه — با محوریت مدیریت زمان}
\textbf{فاز ۱: انتظار اجباری + حفظ سرمایه (ماه‌های ۱–۱۵)}
\begin{itemize}[rightmargin=1.5cm]
\item \textbf{هدف اصلی}: رسیدن به ۲ سال سابقه عملیاتی (۱۵ ماه باقی‌مانده)
\item \textbf{هدف ثانوی}: حفظ حقوق سهام بالای \USD{30M} در طول این دوره
\item \textbf{استراتژی مدیریت نرخ مصرف (\lat{Burn Rate Management})}:
  \begin{enumerate}[rightmargin=1.5cm]
  \item کنترل هزینه‌های عملیاتی غیرضروری
  \item تمرکز منابع بر نقاط عطف علمی (\lat{scientific milestones})
  \item در صورت نیاز: دور تأمین مالی پلی (\lat{bridge round}) کوچک (\USD{5M}) برای پوشش دوره انتظار
  \end{enumerate}
\end{itemize}

\textbf{فاز ۲: آماده‌سازی \lat{IPO} (ماه‌های ۱۶–۱۸)}
\begin{itemize}[rightmargin=1.5cm]
\item تکمیل حسابرسی‌های مالی ۳ ساله (شروع از ماه ۱۲)
\item ایجاد ساختارهای حاکمیت شرکتی (\lat{corporate governance})
\item آماده‌سازی مستندات ثبت \lat{S-1}
\end{itemize}

\textbf{فاز ۳: اجرای \lat{IPO} (ماه‌های ۱۹–۲۴)}
\begin{itemize}[rightmargin=1.5cm]
\item ثبت و تأیید \lat{SEC}
\item اجرای جلسات معرفی به سرمایه‌گذاران (\lat{roadshow})
\item قیمت‌گذاری نهایی و عرضه اولیه
\end{itemize}
\end{processbox}

\subsection{خلاصه تحلیلی و مدیریت ریسک}

\begin{summarybox}{نتیجه‌گیری سناریو و چارچوب زمانی}
\begin{tabularx}{\textwidth}{>{\bfseries\color{headerColor}}l X C{2cm}}
\textbf{فاصله اصلی شناسایی‌شده} & فقط زمانی (۱۵ ماه تا ۲ سال سابقه) & — \\[4pt]
\textbf{راه‌حل عملیاتی پیشنهادی} & صبر اجباری + مدیریت نرخ مصرف سرمایه & — \\[4pt]
\textbf{چارچوب زمانی کل} & ۲۴ ماه (۱۵ ماه انتظار + ۹ ماه فرآیند \lat{IPO}) & \textbf{۲۴ ماه} \\[4pt]
\textbf{ریسک‌های اصلی اجرایی} & نرخ مصرف سرمایه بیش از حد قبل از \lat{IPO}، تأخیر در نقاط عطف علمی & سطح بالا \\[4pt]
\textbf{شاخص‌های نظارتی کلیدی} & حقوق سهام ≥ \USD{30M}، نرخ مصرف ماهانه ≤ \USD{300K}، حفظ تیم کلیدی & — \\
\end{tabularx}
\end{summarybox}

\returnscenarios

% SCENARIO 3
\clearpage
\section{سناریو ۳: یونیکورن \lat{AI} در مرحله \lat{Pre-IPO}}

\begin{examplebox}{پروفایل شرکت و موقعیت بازار}
\textbf{نوع کسب‌وکار:} پلتفرم \lat{AI/ML} برای \lat{enterprise}\\
\textbf{سن شرکت:} ۵ سال\\
\textbf{موقعیت مالی:} رشد فوق‌العاده سریع، زیان‌ده ولی ارزش‌گذاری بالا\\
\textbf{مسیر هدف:} استاندارد ارزش بازار (\lat{Market Value Standard})\\
\textbf{منطق انتخاب مسیر:} ارزش‌گذاری بسیار بالا، زیان‌دهی موقتی، مقیاس بزرگ
\end{examplebox}

\subsection{موقعیت فعلی (نقطه شروع تحلیلی)}

\begin{table}[H]
\centering
\caption{سناریو ۳ — وضعیت مالی و عملیاتی فعلی}
\label{tab:scenario3_current}
\begin{threeparttable}
\begin{tabularx}{\textwidth}{H{4.5cm} X C{2.5cm}}
\toprule
\rowcolor{tableHeader}
\textcolor{white}{\textbf{شاخص کلیدی}} & \textcolor{white}{\textbf{توضیح تحلیلی و اهمیت استراتژیک}} & \textcolor{white}{\textbf{مقدار فعلی}} \\
\midrule
\rowcolor{tableRow1}
درآمد سالانه تکراری (\lat{ARR}) & درآمد مکرر سالانه (رشد ۲۰۰٪ سالانه) & \USD{45M} \\
\rowcolor{conceptBox}
زیان خالص (\lat{Net Loss}) & زیان عملیاتی به دلیل سرمایه‌گذاری سنگین در رشد & \lat{-}\USD{25M} \\
\rowcolor{tableRow1}
حقوق سهام (\lat{Equity}) & منفی به دلیل زیان انباشه از دوره‌های رشد & \lat{-}\USD{8M} \\
\rowcolor{requirementBox}
تأمین مالی تجمعی & کل سرمایه جذب‌شده تا دور \lat{Series D} & \USD{180M} \\
\rowcolor{tableRow1}
ارزش‌گذاری آخرین دور & ارزش‌گذاری در دور \lat{Series D} (۴ ماه پیش) & \USD{900M} \\
\bottomrule
\end{tabularx}
\end{threeparttable}
\end{table}

\subsection{تحلیل فاصله تا استاندارد ارزش بازار}

\begin{table}[H]
\centering
\caption{سناریو ۳ — تحلیل فاصله کمی تا الزامات مسیر ارزش بازار}
\label{tab:scenario3_gap}
\begin{threeparttable}
\begin{tabularx}{\textwidth}{H{4cm} C{2.2cm} C{2.2cm} X}
\toprule
\rowcolor{tableHeader}
\textcolor{white}{\textbf{الزام}} & \textcolor{white}{\textbf{آستانه}} & \textcolor{white}{\textbf{وضعیت فعلی}} & \textcolor{white}{\textbf{تحلیل فاصله و اولویت}} \\
\midrule
\rowcolor{pathwayMarket}
ارزش بازار کل اوراق (\lat{MVLS}) & \USD{75M} >= & \USD{350M}\tnote{a} & \textbf{برآورده شده} — فاصله: صفر \\
\rowcolor{tableRow2}
ارزش بازار سهام عمومی & \USD{20M} >= & \USD{175M}\tnote{b} & \textbf{برآورده شده} — فاصله: صفر \\
\rowcolor{pathwayMarket}
بازارسازهای متعهد & \lat{4} >= & تأمین خودکار & \textbf{برآورده شده} — فاصله: صفر \\
\bottomrule
\end{tabularx}
\begin{tablenotes}\small
\item[a] فرض: \lat{IPO} با ارزش‌گذاری محافظه‌کارانه \USD{700M} (\lat{down-round} از \USD{900M})
\item[b] فرض: \lat{float} ۵۰٪ → \USD{350M} ارزش بازار سهام عمومی
\end{tablenotes}
\end{threeparttable}
\end{table}

\subsection{تحلیل موقعیت و مسیر عملیاتی}

\begin{examplebox}{موقعیت ایده‌آل: بدون فاصله ساختاری}
این شرکت \textbf{هیچ فاصله ساختاری} برای مسیر ارزش بازار ندارد. ارزش‌گذاری \USD{900M} به‌خودی‌خود کافی است، حتی با وجود حقوق سهام منفی و زیان‌دهی. این نشان‌دهنده قدرت مسیر ارزش بازار برای شرکت‌های پررشد است.
\end{examplebox}

\begin{processbox}{برنامه اجرایی ۱۵ ماهه — بهینه‌سازی برای \lat{IPO}}
\textbf{فاز ۱: بهینه‌سازی اقتصاد واحد (ماه‌های ۱–۶)}
\begin{itemize}[rightmargin=1.5cm]
\item \textbf{هدف}: کاهش نرخ مصرف سرمایه (\lat{burn rate}) و بهبود معیارهای واحد
\item \textbf{دلیل استراتژیک}: بهبود داستان سرمایه‌گذاری (\lat{investment story}) برای \lat{IPO}
\item \textbf{اقدامات کلیدی}:
  \begin{enumerate}[rightmargin=1.5cm]
  \item افزایش حاشیه سود ناخالص (\lat{gross margin})
  \item نشان دادن مسیر واضح به سمت سودآوری
  \item تقویت معیارهای نگهداری مشتری (\lat{retention metrics})
  \end{enumerate}
\end{itemize}

\textbf{فاز ۲: فرآیند \lat{IPO} (ماه‌های ۷–۱۵)}
\begin{itemize}[rightmargin=1.5cm]
\item تکمیل حسابرسی‌های مالی ۳ ساله
\item ساختاردهی حاکمیت شرکتی
\item انتخاب بانک‌های سرمایه‌گذاری (\lat{bulge bracket underwriters})
\item ثبت \lat{S-1} و فرآیند بررسی \lat{SEC}
\item جلسات معرفی و قیمت‌گذاری نهایی
\end{itemize}
\end{processbox}

\subsection{خلاصه تحلیلی و مدیریت ریسک}

\begin{summarybox}{نتیجه‌گیری سناریو و چارچوب زمانی}
\begin{tabularx}{\textwidth}{>{\bfseries\color{headerColor}}l X C{2cm}}
\textbf{فاصله اصلی شناسایی‌شده} & هیچ — ارزش‌گذاری کافی است & — \\[4pt]
\textbf{راه‌حل عملیاتی پیشنهادی} & بهینه‌سازی اختیاری برای بهبود داستان \lat{IPO} & — \\[4pt]
\textbf{چارچوب زمانی کل} & ۱۵ ماه (۶ ماه بهینه‌سازی + ۹ ماه \lat{IPO}) & \textbf{۱۵ ماه} \\[4pt]
\textbf{ریسک‌های اصلی اجرایی} & رکود بازار که ارزش‌گذاری را کاهش دهد، تغییر در تمایل بازار به فناوری & سطح پایین \\[4pt]
\textbf{شاخص‌های موفقیت کلیدی} & حفظ ارزش‌گذاری بالای \USD{700M}، بهبود معیارهای واحد & — \\
\end{tabularx}
\end{summarybox}

\returnscenarios

% SCENARIO 4
\clearpage
\section{سناریو ۴: شرکت تولیدی مستقر}

\begin{examplebox}{پروفایل شرکت و موقعیت بازار}
\textbf{نوع کسب‌وکار:} تولیدکننده قطعات صنعتی\\
\textbf{سن شرکت:} ۱۵ سال (\lat{family-owned})\\
\textbf{موقعیت مالی:} سودآور و پایدار، می‌خواهد برای رشد عمومی شود\\
\textbf{مسیر هدف:} استاندارد دارایی/درآمد (\lat{Assets/Revenue Standard})\\
\textbf{منطق انتخاب مسیر:} دارایی فیزیکی بالا، درآمد پایدار، سودآوری اثبات‌شده
\end{examplebox}

\subsection{موقعیت فعلی (نقطه شروع تحلیلی)}

\begin{table}[H]
\centering
\caption{سناریو ۴ — وضعیت مالی و عملیاتی فعلی}
\label{tab:scenario4_current}
\begin{threeparttable}
\begin{tabularx}{\textwidth}{H{4.5cm} X C{2.5cm}}
\toprule
\rowcolor{tableHeader}
\textcolor{white}{\textbf{شاخص کلیدی}} & \textcolor{white}{\textbf{توضیح تحلیلی و اهمیت استراتژیک}} & \textcolor{white}{\textbf{مقدار فعلی}} \\
\midrule
\rowcolor{tableRow1}
درآمد سالانه & فروش پایدار با رشد متوسط & \USD{78M} \\
\rowcolor{conceptBox}
مجموع دارایی‌ها & تجهیزات، موجودی، املاک و مستغلات & \USD{82M} \\
\rowcolor{tableRow1}
سود خالص & \lat{EBITDA} با حاشیه سود ۷٪ & \USD{5.5M} \\
\rowcolor{requirementBox}
حقوق سهام & دارایی‌ها منهای بدهی‌ها & \USD{28M} \\
\rowcolor{tableRow1}
بدهی & سطح معمول برای شرکت تولیدی & \USD{54M} \\
\bottomrule
\end{tabularx}
\end{threeparttable}
\end{table}

\subsection{تحلیل فاصله تا استاندارد دارایی/درآمد}

\begin{table}[H]
\centering
\caption{سناریو ۴ — تحلیل فاصله کمی تا الزامات مسیر دارایی/درآمد}
\label{tab:scenario4_gap}
\begin{threeparttable}
\begin{tabularx}{\textwidth}{H{4cm} C{2.2cm} C{2.2cm} X}
\toprule
\rowcolor{tableHeader}
\textcolor{white}{\textbf{الزام}} & \textcolor{white}{\textbf{آستانه}} & \textcolor{white}{\textbf{وضعیت فعلی}} & \textcolor{white}{\textbf{تحلیل فاصله و اولویت}} \\
\midrule
\rowcolor{pathwayAssets}
مجموع دارایی‌ها & \USD{75M} >= & \USD{82M} & \textbf{برآورده شده} — فاصله: صفر \\
\rowcolor{tableRow2}
مجموع درآمدها & \USD{75M} >= & \USD{78M} & \textbf{برآورده شده} — فاصله: صفر \\
\rowcolor{pathwayAssets}
ارزش بازار سهام عمومی & \USD{20M} >= & \USD{30M}\tnote{a} & \textbf{برآورده شده} — فاصله: صفر \\
\rowcolor{tableRow2}
بازارسازهای متعهد & \lat{4} >= & تأمین خودکار & \textbf{برآورده شده} — فاصله: صفر \\
\bottomrule
\end{tabularx}
\begin{tablenotes}\small
\item[a] فرض: \lat{IPO} با ضریب \lat{P/E} محافظه‌کارانه \lat{10x} → ارزش‌گذاری \USD{55M}، \lat{float} ۵۰٪ → \USD{27.5M}
\end{tablenotes}
\end{threeparttable}
\end{table}

\subsection{تحلیل موقعیت و مسیر عملیاتی}

\begin{examplebox}{موقعیت ایده‌آل: آماده فوری}
این شرکت \textbf{هیچ فاصله ساختاری} ندارد. دارایی و درآمد فعلی کافی است. تنها کار باقی‌مانده، آماده‌سازی اداری و قانونی \lat{IPO} است.
\end{examplebox}

\begin{processbox}{برنامه اجرایی ۱۲ ماهه — تمرکز بر فرآیند \lat{IPO}}
\textbf{فاز ۱: تبدیل به شرکت عمومی (ماه‌های ۱–۳)}
\begin{itemize}[rightmargin=1.5cm]
\item \textbf{حسابرسی تاریخی}: تکمیل حسابرسی‌های \lat{GAAP} برای ۳ سال
\item \textbf{حاکمیت شرکتی}: ایجاد هیئت‌مدیره مستقل
\item \textbf{انطباق \lat{SOX}}: شروع آماده‌سازی کنترل‌های داخلی
\end{itemize}

\textbf{فاز ۲: انتخاب شرکای \lat{IPO} (ماه‌های ۴–۶)}
\begin{itemize}[rightmargin=1.5cm]
\item انتخاب بانک سرمایه‌گذاری (\lat{investment banker})
\item انتخاب وکلای اوراق بهادار (\lat{securities lawyers})
\item تشکیل کمیته حسابرسی (\lat{audit committee})
\end{itemize}

\textbf{فاز ۳: اجرای \lat{IPO} (ماه‌های ۷–۱۲)}
\begin{itemize}[rightmargin=1.5cm]
\item آماده‌سازی و ثبت \lat{S-1}
\item بررسی \lat{SEC} و اصلاحات
\item جلسات معرفی
\item قیمت‌گذاری و فهرست شدن
\end{itemize}
\end{processbox}

\subsection{خلاصه تحلیلی و مدیریت ریسک}

\begin{summarybox}{نتیجه‌گیری سناریو و چارچوب زمانی}
\begin{tabularx}{\textwidth}{>{\bfseries\color{headerColor}}l X C{2cm}}
\textbf{فاصله اصلی شناسایی‌شده} & هیچ — همه معیارها برآورده شده & — \\[4pt]
\textbf{راه‌حل عملیاتی پیشنهادی} & فقط فرآیند اداری \lat{IPO} & — \\[4pt]
\textbf{چارچوب زمانی کل} & ۱۲ ماه (۳ ماه آماده‌سازی + ۹ ماه \lat{IPO}) & \textbf{۱۲ ماه} \\[4pt]
\textbf{ریسک‌های اصلی اجرایی} & کم — شرکت مستقر و سودآور & بسیار پایین \\[4pt]
\textbf{شاخص‌های موفقیت کلیدی} & تکمیل به موقع فرآیندهای نظارتی، حفظ عملکرد مالی & — \\
\end{tabularx}
\end{summarybox}

\returnscenarios

% COMPARATIVE SUMMARY
\clearpage
\section{مقایسه جامع چهار سناریو}

\begin{table}[H]
\centering
\caption{خلاصه مقایسه‌ای سناریوها}
\label{tab:scenarios_comparison}
\begin{threeparttable}
\begin{tabularx}{\textwidth}{H{2.5cm} P{2cm} P{2.5cm} P{2.5cm} C{1.5cm}}
\toprule
\rowcolor{tableHeader}
\textcolor{white}{\textbf{سناریو}} & \textcolor{white}{\textbf{مسیر}} & \textcolor{white}{\textbf{فاصله اصلی}} & \textcolor{white}{\textbf{راه‌حل}} & \textcolor{white}{\textbf{زمان}} \\
\midrule
\rowcolor{pathwayIncome}
۱. \lat{SaaS} & درآمدی & \USD{6M} سرمایه + \USD{575K} سود & رشد + تأمین مالی & ۱۸ ماه \\
\rowcolor{pathwayEquity}
۲. \lat{Biotech} & سرمایه & ۱۵ ماه سابقه & صبر اجباری & ۲۴ ماه \\
\rowcolor{pathwayMarket}
۳. \lat{AI Unicorn} & ارزش بازار & هیچ & بهینه‌سازی اختیاری & ۱۵ ماه \\
\rowcolor{pathwayAssets}
۴. تولیدی & دارایی/درآمد & هیچ & فرآیند \lat{IPO} & ۱۲ ماه \\
\bottomrule
\end{tabularx}
\end{threeparttable}
\end{table}

\begin{summarybox}{نتیجه‌گیری کلیدی}
\begin{itemize}[rightmargin=1.5cm]
\item \textbf{سریع‌ترین مسیر}: شرکت‌های تولیدی مستقر با دارایی/درآمد بالا (۱۲ ماه)
\item \textbf{کندترین مسیر}: استارتاپ‌های جوان که به ۲ سال سابقه نیاز دارند (۲۴ ماه)
\item \textbf{مسیر قابل کنترل}: رشد به سودآوری برای \lat{SaaS} (۱۸ ماه با اجرای خوب)
\item \textbf{مسیر بدون مانع}: یونیکورن‌ها با ارزش‌گذاری بالا (۱۵ ماه، بدون فاصله ساختاری)
\end{itemize}
\end{summarybox}

\returnhome

% ===================== STRATEGIC INSIGHTS =====================
\clearpage
\chapter{نکات استراتژیک}

\section{انتخاب مسیر بهینه}

\begin{principlebox}{اصول تصمیم‌گیری}
انتخاب مسیر باید بر اساس \textbf{نقاط قوت فعلی} شرکت باشد، نه آرزوها:

\begin{itemize}[rightmargin=1.5cm]
\item اگر \textbf{سودآور} هستید → مسیر درآمدی (سریع‌ترین برای شرکت‌های سودآور)
\item اگر \textbf{سرمایه قوی} ولی جوان هستید → مسیر سرمایه (اما باید ۲ سال صبر کنید)
\item اگر \textbf{ارزش‌گذاری بالا} دارید → مسیر ارزش بازار (راه میانبر برای یونیکورن‌ها)
\item اگر \textbf{دارایی و درآمد بالا} دارید → مسیر دارایی/درآمد (ایده‌آل برای صنایع)
\end{itemize}
\end{principlebox}
\clearpage
\section{تله‌های رایج و پیشگیری}

\begin{warningbox}{اشتباهات متداول و راه‌حل‌ها}
\begin{enumerate}[rightmargin=1.5cm]
\item \textbf{\lat{Premature IPO}}: رفتن به بازار قبل از آماده بودن الزامات
  \begin{itemize}
  \item \textbf{خطر}: ریزش قیمت در روزهای اول، آسیب به اعتبار
  \item \textbf{پیشگیری}: تحلیل دقیق فاصله‌ها قبل از شروع فرآیند
  \end{itemize}

\item \textbf{انتخاب مسیر اشتباه}: انتخاب مسیری که با نقاط قوت شرکت همخوانی ندارد
  \begin{itemize}
  \item \textbf{مثال}: یک شرکت بیوتکنولوژی پیش‌درآمد تلاش برای مسیر درآمدی (غیرممکن)
  \item \textbf{پیشگیری}: تحلیل واقع‌بینانه از وضعیت فعلی شرکت
  \end{itemize}

\item \textbf{فرض ارزش‌گذاری غیرواقعی}: بیش‌ازحد تخمین زدن تمایل بازار
  \begin{itemize}
  \item \textbf{مشکل}: ارزش بازار سهام عمومی بر اساس ارزش‌گذاری \lat{IPO} واقعی است، نه آخرین دور خصوصی
  \item \textbf{پیشگیری}: استفاده از ارزیابی‌های محافظه‌کارانه
  \end{itemize}

\item \textbf{نادیده گرفتن الزامات زمانی}: فراموش کردن ۲ سال سابقه برای استاندارد سرمایه
  \begin{itemize}
  \item \textbf{محدودیت}: این یک مانع سخت است — نمی‌توان با پول حل کرد
  \item \textbf{پیشگیری}: برنامه‌ریزی بلندمدت از ابتدا
  \end{itemize}
\end{enumerate}
\end{warningbox}

\section{بهینه‌سازی زمان‌بندی}

\begin{processbox}{استراتژی‌های تسریع زمان تا \lat{IPO}}
برای کوتاه کردن زمان تا \lat{IPO}:

\textbf{اقدامات موازی‌سازی‌شده:}
\begin{itemize}[rightmargin=1.5cm]
\item شروع حسابرسی‌های تاریخی همزمان با رفع فاصله‌ها
\item ساختاردهی حاکمیت شرکتی قبل از رسیدن به آستانه‌ها
\item انتخاب مشاوران \lat{IPO} در مراحل اولیه
\end{itemize}

\textbf{اقدامات پیشگیرانه:}
\begin{itemize}[rightmargin=1.5cm]
\item ثبت مالیات و حسابداری طبق \lat{GAAP/IFRS} از ابتدا
\item ایجاد کنترل‌های داخلی قبل از الزام \lat{SOX}
\item داشتن مدیران مستقل از قبل
\end{itemize}
\end{processbox}

\returnhome

% ===================== APPENDIX: GLOSSARY =====================
\clearpage
\chapter*{واژه‌نامه تخصصی}
\addcontentsline{toc}{chapter}{واژه‌نامه}

\begin{table}[H]
\centering
\caption{اصطلاحات کلیدی مالی و حقوقی}
\label{tab:glossary}
\begin{threeparttable}
\begin{tabularx}{\textwidth}{P{4cm} X}
\toprule
\rowcolor{tableHeader}
\textcolor{white}{\textbf{اصطلاح}} & \textcolor{white}{\textbf{تعریف عملیاتی و کاربرد}} \\
\midrule
\rowcolor{conceptBox}
\lat{GAAP} & استانداردهای حسابداری پذیرفته‌شده عمومی آمریکا — مبنای گزارش‌گری مالی برای شرکت‌های آمریکایی \\
\rowcolor{tableRow1}
\lat{IFRS} & استانداردهای بین‌المللی گزارشگری مالی — استاندارد جهانی گزارش‌گری مالی \\
\rowcolor{requirementBox}
\lat{MVLS} & ارزش بازار کل اوراق قابل معامله — معیار اصلی برای مسیر ارزش بازار \\
\rowcolor{tableRow1}
\lat{MVUPHS} & ارزش بازار سهام عمومی بدون محدودیت — معیار نقدشوندگی اولیه \\
\rowcolor{conceptBox}
\lat{PCAOB} & هیئت نظارت بر حسابداری شرکت‌های عمومی — نهاد نظارتی حسابرسان در آمریکا \\
\rowcolor{tableRow1}
\lat{SOX} & قانون ساربنز-آکسلی — الزامات کنترل‌های داخلی برای شرکت‌های عمومی \\
\rowcolor{requirementBox}
\lat{S-1} & فرم ثبت اولیه نزد \lat{SEC} برای \lat{IPO} — سند اصلی ثبت عرضه اولیه \\
\rowcolor{tableRow1}
\lat{Float} & درصد سهامی که در اختیار عموم قرار می‌گیرد — معمولاً ۳۰–۵۰٪ در \lat{IPO} \\
\rowcolor{conceptBox}
\lat{Lock-up} & دوره‌ای که سهام دارندگان داخلی قابل فروش نیست — معمولاً ۱۸۰ روز پس از \lat{IPO} \\
\rowcolor{tableRow1}
\lat{Underwriter} & بانک سرمایه‌گذاری مدیریت‌کننده \lat{IPO} — مسئول قیمت‌گذاری و توزیع سهام \\
\rowcolor{requirementBox}
\lat{Roadshow} & جلسات معرفی شرکت به سرمایه‌گذاران نهادی قبل از \lat{IPO} \\
\rowcolor{tableRow1}
\lat{Greenshoe} & اختیار خرید اضافی ۱۵٪ برای بانک‌های سرمایه‌گذاری در صورت تقاضای بالا \\
\bottomrule
\end{tabularx}
\end{threeparttable}
\end{table}

\returnhome

\begin{figure}
    \centering
    \includegraphics[width=1\linewidth]{IPO_USABLE_English.png}
\end{figure}

% ===================== FINAL PAGE =====================
\clearpage
\thispagestyle{empty}

\begin{center}
\vspace*{3cm}

\begin{tikzpicture}[scale=1.5, transform shape]
    \node[pathway1, minimum width=3.5cm] (p1) at (0,2) {\Large\textbf{درآمدی}\\[-2pt]\small \USD{15M} + \USD{1M} سود};
    \node[pathway2, minimum width=3.5cm] (p2) at (3.5,0) {\Large\textbf{سرمایه}\\[-2pt]\small \USD{30M} + ۲ سال};
    \node[pathway3, minimum width=3.5cm] (p3) at (0,-2) {\Large\textbf{ارزش بازار}\\[-2pt]\small \USD{75M} ارزش‌گذاری};
    \node[pathway4, minimum width=3.5cm] (p4) at (-3.5,0) {\Large\textbf{دارایی/درآمد}\\[-2pt]\small \USD{75M} هر یک};
    
    \node[draw=headerColor, fill=white, circle, minimum size=2.5cm, line width=2pt] at (0,0) {\LARGE\bfseries\color{headerColor} \lat{Nasdaq}};
    
    \draw[arrow, headerColor] (p1) -- (0,0.9);
    \draw[arrow, headerColor] (p2) -- (1.6,0);
    \draw[arrow, headerColor] (p3) -- (0,-0.9);
    \draw[arrow, headerColor] (p4) -- (-1.6,0);
\end{tikzpicture}

\vspace{2.5cm}

{\color{headerColor}\rule{0.7\textwidth}{0.75pt}}

\vspace{0.8cm}

{\LARGE\bfseries\color{headerColor}پایان سند}

\vspace{0.3cm}

{\Large\bfseries مسیرهای ورود به \lat{Nasdaq Global Market}}

\vspace{0.3cm}

{\large تحلیل چهار مسیر استاندارد — نسخه ۱٫۰}

\vspace{0.8cm}

{\color{gray} تاریخ انتشار: ۱۰ بهمن ۱۴۰۴ (\lat{30 January 2026})}

\vspace{0.3cm}

{\color{gray}\small منبع اولیه: \lat{Nasdaq Rule 5405} — \lat{Global Market Initial Listing Standards}}

\end{center}

\end{document}